\chapter*{Glossário do Capítulo 5}
\addcontentsline{toc}{chapter}{Glossário do Capítulo 5}

\begin{description}[
  font=\normalfont\bfseries,
  style=multiline,
  leftmargin=5.8cm,
  labelwidth=5.1cm,
  labelsep=0.35cm,
  itemsep=0.55\baselineskip
]
  \item[CMake] Sistema de configuração e geração de build usado pelo ESP-IDF para descrever componentes, dependências e alvos de compilação.
  \item[CTR\_DRBG] Mecanismo determinístico de geração de números pseudoaleatórios baseado em cifra de bloco, usado para apoiar geração de chaves e assinaturas.
  \item[ESP-IDF] Sigla de \textit{Espressif IoT Development Framework}, framework oficial de desenvolvimento para dispositivos da família ESP32.
  \item[ESP32] Família de sistemas em chip da Espressif voltada a aplicações embarcadas conectadas, com suporte a Wi-Fi, periféricos seriais e execução local de firmware.
  \item[Event Loop] Mecanismo central do ESP-IDF para entrega de eventos do sistema, como mudança de estado da rede ou desconexão do Wi-Fi.
  \item[FreeRTOS] Sistema operacional de tempo real usado como base de agendamento, tarefas, filas e sincronização no ESP-IDF.
  \item[Firmware] Camada de software executada diretamente no dispositivo embarcado, responsável por controlar periféricos, rede e lógica local.
  \item[GPIO] Sigla de \textit{General-Purpose Input/Output}, pinos programáveis usados para entrada e saída digital.
  \item[Mbed TLS] Biblioteca criptográfica integrada ao ESP-IDF, usada para hash, TLS, geração de chaves e assinaturas digitais.
  \item[NVS] Sigla de \textit{Non-Volatile Storage}, mecanismo de armazenamento persistente usado pelo ESP-IDF para dados locais e apoio a componentes como Wi-Fi.
  \item[PlatformIO] Ecossistema de build e automação que integra toolchain, gerenciamento de plataforma, upload e monitor serial para projetos embarcados.
  \item[SDK] Sigla de \textit{Software Development Kit}, conjunto de bibliotecas, ferramentas e documentação necessário para desenvolver sobre uma plataforma.
  \item[SoC] Sigla de \textit{System on Chip}, circuito integrado que reúne CPU, memória e periféricos em um único chip.
  \item[SPI Host] Controlador SPI do lado mestre no ESP32, responsável por iniciar e coordenar transações no barramento.
  \item[Station Mode] Modo em que o ESP32 atua como cliente Wi-Fi e se conecta a um ponto de acesso já existente.
  \item[Task] Unidade de execução agendada pelo FreeRTOS, comparável a uma thread em sistemas embarcados.
\end{description}
